% 2011 年山东卷理科第 11 题：两个全等矩形的正（主）视图与俯视图。
% 来源 PDF：普通高考/2011/2011山东理.pdf，SHA-256 d0f409fd3fa82b12c210d587c8412d4c446946a201f87d124e2fc17c5a50db16。
% 原页：第 1 页（A3 三栏）第 2 栏中部，书页码 689。
% 重绘依据：照原页——左、右两个同样大小（长宽比约 5:3）的矩形并排，图上无尺寸标注，
%   仅在下方分别标「正（主）视图」「俯视图」；两矩形外框为各自视图的轮廓，不画投影连线。
\begin{tikzpicture}[scale=1.0, semithick]
  \draw (0,0) rectangle (2.4,1.45);
  \node[below] at (1.2,-0.12) {正（主）视图};
  \draw (3.3,0) rectangle (5.7,1.45);
  \node[below] at (4.5,-0.12) {俯视图};
\end{tikzpicture}
